\uuid{VT5z}
\titre{Pgcd}
\auteur{bodin ; croizat ; sacre}
\organisation{exo7}
\licence{CC BY-NC-SA 4.0}
\datecreate{2022-03-08}
\niveau{L1}
\module{Algèbre}
\chapitre{Arithmétique dans Z}
\sousChapitre{Pgcd, ppcm, algorithme d'Euclide}
\difficulte{3}
\theme{pgcd}
\exo7qcmid{hilisit.arithmetique.012}

\contenu{
\begin{question}
\qclassification{103.03}
    Soit $a_1$ un entier dont le reste dans la division euclidienne par $5$ est $r_1 = 2$. Soit $a_2$ un entier dont le reste dans la division euclidienne par $5$ est $r_2 = 3$. Quelles sont alors les affirmations vraies ?
    \begin{answers}
        \good{Le reste de la division euclidienne de $a_1 + a_2$ par $5$ est $0$.}
        \bad{Le reste de la division euclidienne de $a_1 + a_2$ par $5$ est $5$.}
        \good{Le reste de la division euclidienne de $2a_1 + 2a_2$ par $5$ est $0$.}
        \good{L'écriture décimale de $2a_1 + 2a_2$ finit par le chiffre $0$.}
    \end{answers}
    \begin{explanations}
    Les divisions euclidiennes par $5$ nous donnent : $a_1 = 5k_1 + 2$ et $a_2 = 5k_2 + 3$. On a ainsi :
    $$ a_1 + a_2 = 5(k_1+k_2) + 5 = 5(k_1+k_2+1) + 0 $$
    La dernière écriture correspond bien à la division euclidienne de $a_1 + a_2$ par $5$ car le reste (c'est $0$) est bien compris entre $0$ et $4$.\\
    De même, on calcule :
    $$ 2a_1 + 2a_2 = 2 \times 5(k_1+k_2+1) = 5 (2k_1 + 2k_2 + 2) = 10 (k_1 + k_2 + 1) $$
    Aussi $2a_1 + 2a_2$ est un entier divisible par $5$ et par $10$, donc son écriture décimale se termine par $0$.
    \end{explanations}
\end{question}
}
